% critique.tex — f171r, vv. 721–763 : commentaire critique.
% Articles textuels tirés du commentaire de l'auteur (latex/Aristophane-Thesmo Harmattan 2023 vv 642-fin-bis.tex),
% séparés du commentaire de la pièce (commentaire.tex) ; texte et notes conservés. Détail : repartition.md.
\cc{721}{ξὺν ἀθέοις ἔργοις}{(Voir au commentaire de la pièce, vv. 721–5, la remarque de Jean Taillardat sur ξὺν ἀδίκοις ἔργοις, v. 716.) Fort de cette remarque, je [Bibi] n'ai pas hésité à introduire ce ξύν dans le texte  au vers 721, qui avait un problème de mètre : [ξὺν] ἀθέοις ἔργοις. C'est aussi \italique{gratia metri} qu'Enger proposait en son temps l'ajout d' [ἐπ'] ἀθέοις ἔργοις, conjecture acceptée par Lætitia (\italique{The Songs of Aristophanes}, p. 427). N.D.É.}
\cc{724}{τάχα δὲ \{σε\} μεταβαλοῦσ᾽}{Quant à σε de τάχα δέ σε μεταβαλοῦσα, κτλ. (cod. R), Jean Taillardat le considère comme une glose qui expliquait κακόν ἑτερότροπον et qui s'est glissée dans le texte même. (Voir le commentaire de la pièce, vv. 721–5.)}
\cc{728}{ἴωμεν ἐπὶ τὰς κληματίδας.}{Blaydes, jugeant l'article de trop, propose : ἴωμεν οὖν ἐπὶ κληματίδας. Nous répondrons par le texte suivant de Plutarque, où il s'agit, comme ici, d'un sanctuaire où des femmes se trouvent réunies pour la célébration de leurs mystères :\citComm{$[$ ἡ Δαμοκρίτα $]$ ξύλα ταῖς θύραις πολλὰ προσνήσασα ταῦτα δ᾽ ἦν \textbf{εἰς τὴν τῆς ἑορτῆς θυσίαν ὑπ᾽ ἐκείνων παρεσκευασμένα}, πῦρ ἐνῆκε.\\
Damocrita ayant entassé contre les portes une quantité de bois, \italique{que les femmes avaient préparé pour le sacrifice de la fête}, elle y mit le feu.}(Plutarque, \italique{Narrat. amour}., 5, p. 775e ; Note de Willems.)}
\cc{742–66}{}{L'on possède pour ce passage un autre témoin que le \italique{Ravennas} : le premier fragment du \nieme{3839} papyrus d'Oxyrhynque\nc{\italique{The Oxyrhynchus Papyri}, vol. LVI (1989), n°3839, fr. 1.} (\PsiDeux \  dans l'apparat critique). Ce morceau de papyrus est bien antérieur au manuscrit de Ravenne : il est daté d'environ 200 ap. J.-C. Il présente objectivement de meilleures leçons en bien des points. Aussi, aux endroits où ces deux témoins différaient et que le bon choix à faire n'était pas évident, j'ai privilégié la version du papyrus. Je reproduis en annexe ce papyrus ainsi que la lecture qu'en fit Walter E. H. Coockle pour la collection \italique{The Oxyrhynchus Papyri}.}
\cc{743}{τρικότυλον; ἢ πῶς. --- εἰπέ μοι·}{Comme annoncé, je suis le papyrus d'Oxyrhynque contre le \italique{Ravennas} ; on peut y déchiffrer :\citComm{\_τρικο\douteux{τ}\douteux{υ}\douteux{λ}\douteux{ο}ν·η\douteux{π}\douteux{ω}\douteux{c} \ \ ειπ\douteux{ε}\douteux{μ}οι\douteux{τ}[}Le point médian est, dans ce papyrus, un signe de ponctuation interrogatif. Quant à l'espace blanc de deux lettres entre \douteux{π}\douteux{ω}\douteux{c} et ειπ\douteux{ε}\douteux{μ}οι, il indique que le changement d'interlocuteur a lieu à cet endroit, et non après εἰπέ μοι comme le placent le \italique{Ravennas} et de nombreux éditeurs.
La lecture du papyrus donne donc pour le vers 743 :\citComm{τρικότυλον; ἢ πῶς. --- εἰπέ μοι·τί μ᾽ ἠργάσω;}}
\cc{754}{δός μοι τὸ σφαγεῖον.}{Le papyrus d'Oxyrhynque n°3839 ne comporte pas le μοι. J'ai dit précédemment suivre le papyrus contre le \italique{Ravennas} et je m'y suis tenu. Mais il y a d'excellentes raisons à cela, outre bien-sûr le privilège de l'ancienneté.
D'abord, il n'est rien de plus courant que l'interpolation de petits mots tels que μοι (cf. le vers 1047, où un autre papyrus, daté du \nieme{II} siècle ap. J.-C., ne présente pas non plus le μοι du manuscrit de Ravenne).
Comment ce μοι est-il apparu ? il est vraisemblable que ce soit pour de pseudo-raisons métriques.
Faisons l'hypothèse que d'anciens manuscrits comportaient un \nieme{754} vers ainsi troussé :\citComm{οἴμοι τέκνον δός τὸ σφαγίον Μανία.}Noter la graphie σφαγίον (et non σφαγεῖον) : c'est celle de la première version du \italique{Ravennas} (avant correction en σφαγεῖον, faite d'une autre main s'il faut en croire Coockle, \latin{scholiorum manus, aequalis fere primae manui} selon von Velsen\nc{\italique{Aristophanis Thesmophoriazusae} (1883). Adolphe von Velsen a établi définitivement dans cette édition que le second manuscrit par lequel nous sont connues les \italique{Thesmophories}, le \italique{Monacensis gr. 492} (G dans son apparat critique, \Mon dans le mien), n'était qu'une copie du \nieme{XV} siècle du \italique{Ravennas}, apportant quelques émendations bien vues. Son édition compare exhaustivement le \italique{Ravennas} et le \italique{Monacensis} ; ce qui est d'un grand intérêt pour déterminer si une correction du \italique{Ravennas} est ancienne ou plus récente : il suffit d'aller comparer avec la leçon du manuscrit de Münich. Son apparat critique est le plus précis et permet d'avoir la meilleure idée possible du manuscrit de Ravenne -- avec celui de l'édition Prato ; en partie une reprise du premier mais revu et amélioré en de nombreux points. Pour qui voudrait pousser le scrupule plus loin, et puisque le manuscrit de Ravenne n'a toujours pas été numérisé -- ô scandale ! le manuscrit sans qui nous ne connaîtrions que neuf des onze pièces survivantes d'Aristophane !.. --, Van Leeuwen édita en fac-similé le \italique{Ravennas} : \italique{Comoediae Undecim cum Scholiis : Codex Ravennas 137, 4, A} [$=$ \italique{Rav.} 429 (nouvelle cote)], \italique{phototypice editus} in \italique{Codices Graeci et Latini photographice depicti}, t. IX (1904).}). σφαγίον était donc scandé : $ \cup \cup - $ ; il manquait par conséquent une syllabe longue à placer entre τέκνον et σφαγίον pour que le vers soit d'équerre. L'ajout d'un μοι était la solution la plus simple et à la portée de n'importe quel scribe.\\
Ce sont là des considérations métriques qui devaient passer bien au-dessus de la tête des scribes ayant retranscrit les \italique{Thesmophories} ? Détrompez-vous : la correction même du \italique{Ravennas} -- alors qu'à l'époque les deux graphies se prononçaient de la même façon -- montre bien l'intérêt pour la scansion ; sur le papyrus, Coockle releva une «~\anglais{unexplained ink like an oddly shaped perispomen accent in a second hand above} γε~» (le γε de σφαγειον) ; l'explication me paraît simple : l'annotateur confirmait que la deuxième syllabe de σφαγεῖον est longue. C'est d'ailleurs l'objet d'une glose de Photius (σ 850) :\citComm{σφαγεῖον· [...] περισπᾶται ἡ δευτέρα συλλαβή.}}
\cc{759}{}{Je suis Sommerstein qui donne la parole pour ce vers à Critylla, quoiqu'aucun \italique{paragraphos} indiquant un changement d'interlocuteur ne soit présent sur le manuscrit de Ravenne (le papyrus est inexploitable sur ce point).}
\cc{761}{σοὐξηράσατο}{Je fais mienne la conjecture de Fritzsche : 'ξηράσατο (de ἐξερᾶν) à la place du 'ξῃρήσατο (de ἐξαιρεῖν) présent dans le manuscrit de Ravenne. (Je m'en explique en fin de note.)\\
Puisque le sens premier de la question de Critylla est bien σοι ἐξερᾶν, je reprends la conjecture de Fritzsche (lire σοὐξηράσατο à la place du σοὐξῃρήσατο du \italique{Ravennas}) d'autant que ἐξῃρήσατο serait un aoriste inhabituel\nc{Impossible selon le rigoriste O. Lautensach, p. 128 de sa grammaire \italique{Die Aoriste bei den Attischen Tragikern und Komikern} (1911) : \allemand{Aoriste auf} -ησα\allemand{. Der im Etymologicum Magnum p.777, 53 als "ungriechisch" verworfene und von Lobeck zu Phrynicus p. 716ff. aus der späteren Gräcität nachgewiesene Aor. }ᾕρησα \allemand{ist unmöglich bei Ar. \italique{Thesm.} 761 in} σοὐξῃρήσατο\allemand{ zu dulden, wofür Meineke} σου διεχρήσατο\allemand{, Fritzsche }σου ᾽ξηράσατο\allemand{, Dindorf} σου ᾽ξαφείλετο\allemand{ schreibt, während Lobeck, v. Velsen und Blaydes den Vers für unecht erklären.}}. (Pour le sens, voir le commentaire de la pièce, vv. 760–1.)}
